La tesi ha tracciato l'evoluzione del visual anomaly detection industriale attraverso lo sviluppo dei metodi e la progressione dei dataset di riferimento. Il percorso, che parte dal 2018 fino ai lavori più recenti del 2025, evidenzia come il campo sia in continua crescita.

\medskip

Il \textbf{Capitolo 1} ha definito le basi concettuali necessarie a comprendere questa evoluzione: i principi del machine learning, l'architettura delle reti neurali convoluzionali e il ruolo del transfer learning, alla base dei backbone pre-addestrati su ImageNet utilizzati per estrarre rappresentazioni senza supervisione esplicita sulle anomalie.

\medskip

Il \textbf{Capitolo 2} ha seguito lo spostamento progressivo da approcci generativi di prima generazione (autoencoder e GAN, limitati dalla difficoltà di modellare distribuzioni complesse e dall'instabilità dell'addestramento) verso metodi feature-based che usano backbone pre-addestrati (PaDiM, SPADE, PatchCore) con performance superiori e maggiore stabilità. In fasi successive, i normalizing flow hanno introdotto densità esatte sulla distribuzione delle feature normali; le tecniche di sintesi di anomalie (CutPaste, DRAEM) hanno trasformato il problema in una classificazione binaria; i modelli student-teacher avanzati (EfficientAD) hanno bilanciato accuratezza ed efficienza computazionale. Infine, i vision-language model e i framework multimodali per dati 3D hanno aperto nuovi ambiti di ricerca, pur rivelando limiti ancora significativi: i VLM attuali non competono con i metodi feature-based su benchmark standard, e i backbone 3D esistenti non scalano ai dataset ad alta risoluzione.

\medskip

Il \textbf{Capitolo 3} ha seguito l'evoluzione dei dataset. MVTec AD ha definito il primo benchmark di riferimento, portando ad un primo confronto. I dataset successivi hanno progressivamente ampliato il problema: dalle anomalie logiche di MVTec LOCO AD alle informazioni geometriche di MVTec 3D-AD e Real3D-AD, fino agli scenari large-scale e multi-view di Real-IAD e alle condizioni più realistiche di MVTec AD 2 e RAD. Ogni nuovo benchmark ha evidenziato i limiti dei precedenti, orientando la ricerca verso soluzioni più robuste e generalizzabili.

\medskip

Dall'analisi complessiva emergono alcune osservazioni trasversali. In primo luogo, il rapporto tra le \textit{performance su benchmark controllati} e la \textit{robustezza in scenari reali}: l'ottimizzazione su un singolo dataset non garantisce la generalizzazione. 
In secondo luogo, l'integrazione multimodale (RGB, depth, point cloud, mappe delle normali) rappresenta la direzione migliore per affrontare sempre maggiori categorie di difetti (che possono essere visibili in alcune modalità e non in altre), sviluppando architetture capaci di fondere le informazioni.
In terzo luogo, la disponibilità di dati rimane uno dei vincoli principali: la rarità dei difetti reali e i costi di acquisizione spingono verso approcci few-shot, zero-shot e l'uso di dati sintetici, ma il gap tra simulazione e realtà resta ampio.

\medskip

Le direzioni di ricerca più rilevanti per il futuro convergono verso tre obiettivi: lo sviluppo di modelli capaci di operare su categorie eterogenee senza re-addestramento, la costruzione di architetture 3D scalabili ad alta risoluzione e il passaggio da sistemi di \textit{detection} a sistemi di \textit{comprensione}, capaci di segnalare un'anomalia ma anche di classificarla. 
Il progresso dipenderà dalla disponibilità di benchmark sempre più rappresentativi delle condizioni operative reali.

\medskip

In sintesi, il visual anomaly detection industriale si configura oggi come un campo ben definito nella sua metodologia di base, ma ancora aperto sulle sfide più avanzate. 
Il percorso seguito in questa tesi, dalla detection di graffi su superfici strutturate alla comprensione di assemblaggi complessi in ambienti non controllati, riflette la relazione tra la solidità dei risultati ottenuti e l'ambizione verso sistemi sempre più vicini alle esigenze dell'industria reale.
